\subsection{角的测度}

由第 1 号的定理 1，角的拓扑群 \(\mathfrak{A}\) 同构于 \(\mathbf{T}\) 。\(\mathbf{R}\) 到 \(\mathfrak{A}\) 上的每个严格态射都可由同构 \(z \to \operatorname{Am}(z)\) （\(\mathbf{U}\) 到 \(\mathfrak{A}\) 上的）与 \(\mathbf{R}\) 到 \(\mathbf{U}\) 上的一个严格态射复合而得；若令 \(\Im(x) = \operatorname{Am}(e(x))\) ，则 \(\mathbf{R}\) 到 \(\mathfrak{A}\) 上的每个严格态射都形如 \(x \to \Im(x/a)\) （ \(a \neq 0\) ）。给定一个一劳永逸地固定的实数 \(a > 0\) ，每个角 \(\theta\) 经同态 \(x \to \Im(x/a)\) 对应于实数模 \(a\) 的一个类（即 \(\mathbf{Z}/a\mathbf{Z}\) 的一个元素），它称为 \(\theta\) 相对于底 \(a\) 的测度；按惯用说法，该类中的每个实数也称为 \(\theta\) 的一个测度；角 \(\Im(x/a)\) 称为（相对于底 a）测度为 \(x\) 的角。若 x 是 \(\theta\) 的一个测度而 \(x'\) 是 \(\theta'\) 的一个测度（相对于同一个底），则 \(x + x'\) 是 \(\theta + \theta'\) 的一个测度，且 -x 是 - \(\theta\) 的一个测度。角（相对于底 a）的主测度是其位于区间 {[}0, a{]} 中的那个测度。

底 \(a\) 的选取。我们总限于取 \(a > 1\) 的底。每个 \(a > 1\) 对应一个角 \(\omega = \Im(1/a)\) ，它的主测度是 \(1\) ，称为相对于底 \(a\) 的角的测度单位；反之，每个角 \(\omega \neq 0\) 对应唯一的 \(a > 1\) 使 \(\Im(1/a) = \omega\) ，因此知道角的测度单位就完全确定了底 \(a > 1\) 。

在数值计算中通常取 \(a = 360\) 或 \(a = 400\) ；相应的角的测度单位称为度（ \(a = 360\) ）或百分度（ \(a = 400\) ）。

在分析中，以及在一切不涉及数值计算的数学分支中，普遍使用由条件

\[
\lim _ {x \rightarrow 0} \frac {e (x / a) - 1}{x} = 1
\]

所定义的底 a；这个底记作 \(2\pi\) 。相应的角的测度单位称为弧度，其测度称为弧度测度；利用前面提到的对复数 z 的 \(e^{z}\) 的定义，我们有 \(\boldsymbol{e}(x)=e^{2\pi ix}\) 对一切 \(x\in R\) 成立。

底 a 一经选定，当人们说到一个角时，通常指的就是这个角相对于底 a 的一个测度；只要（在不涉及数值计算时情形总是如此）底 a 自始至终保持固定，并且记住模 a 同余的两个实数对应于同一个角，这种惯用说法并无妨碍。

例如，通常所理解的复数 \(z \neq 0\) 的辐角，是指该角的一个弧度测度，它由依所考虑的问题而定的约定来确定；这些约定一经作出，如此选定的辐角的测度就记作 Am (z)。

